\setcounter{chapter}{0}
\renewcommand{\thetheorem}{\thechapter.\arabic{theorem}}
\setcounter{theorem}{0}

\chapter{The Fibonacci Sequence}

\subsection{Introduction and Generalization}

We begin our study with the classical Fibonacci sequence, a sequence of numbers that naturally emerges in various mathematical and biological contexts.

\begin{definition}[The Fibonacci Sequence]
The \qt{Fibonacci sequence} is an infinite sequence of integers $a_0, a_1, a_2, a_3, \dots, a_n, \dots$ defined by the following recurrence relation:
\begin{equation} \label{eq:standard_fib_recurrence}
\begin{cases}
a_0 = 0 \\
a_1 = 1 \\
a_n = a_{n-1} + a_{n-2} & \text{for } n \ge 2
\end{cases}
\end{equation}
\end{definition}

\begin{remark}
This sequence was introduced to Western mathematics by Leonardo Fibonacci, a medieval mathematician from Pisa (1170--1250), to model the growth of rabbit populations. Computing the first few terms, we easily find:
\[
a_0 = 0, \quad a_1 = 1, \quad a_2 = 1, \quad a_3 = 2, \quad a_4 = 3, \quad a_5 = 5, \quad a_6 = 8, \quad a_7 = 13, \quad a_8 = 21, \dots
\]
\end{remark}

A natural, central question arises: Is there a closed, explicit formula for the general term $a_n$? To answer this, let us briefly look at other well-known mathematical sequences for comparison:
\begin{enumerate}
    \item \textbf{Arithmetic Sequence:} Defined by $a_0 = a$ and $a_n = a_{n-1} + d$ for $n \ge 1$. The explicit formula for the general term is given by $a_n = a + n \cdot d$.
    \item \textbf{Geometric Sequence:} Defined by $a_0 = a$ and $a_n = a_{n-1} \cdot q$ for $n \ge 1$. The explicit formula is given by $a_n = a \cdot q^n$.
\end{enumerate}

Clearly, the recurrence relation for the Fibonacci sequence is fundamentally different. To find a formula for $a_n$, it is highly beneficial to step back and generalize the problem.

\begin{definition}[Generalized Fibonacci Sequence]
Let $a_0, a_1 \in \mathbb{R}$ be given initial real numbers. We define a sequence $\mathcal{A}$ to be a \qt{Fibonacci sequence} if it satisfies the relation:
\begin{equation} \label{eq:generalized_fib_recurrence}
\begin{cases}
\alpha_0 = a_0 \\
\alpha_1 = a_1 \\
\alpha_n = \alpha_{n-1} + \alpha_{n-2} & \text{for } n \ge 2
\end{cases}
\end{equation}
We denote such a sequence uniquely determined by its starting values as $\mathcal{F}_{a_0, a_1}$. Furthermore, we denote the collection (or space) of all such generalized Fibonacci sequences by $\Fib$.
\end{definition}

Notice that our original, classical Fibonacci sequence is precisely $\mathcal{F}_{0, 1}$. 

\subsection{The Algebraic Structure of $\Fib$}

The collection $\Fib$ possesses a very rich internal algebraic structure, which we will exploit to find our explicit formula.

\begin{proposition}[Closure Properties of $\Fib$] \label{prop:fib_closure}
The collection of Fibonacci sequences, $\Fib$, is closed under term-wise addition and scalar multiplication. Specifically,
\begin{enumerate}
    \item If $\mathcal{A}_1, \mathcal{A}_2 \in \Fib$, then $\mathcal{A}_1 + \mathcal{A}_2 \in \Fib$.
    \item If $\mathcal{A} \in \Fib$ and $\alpha \in \mathbb{R}$, then $\alpha \mathcal{A} \in \Fib$.
\end{enumerate}
\end{proposition}

\begin{proof}
For the first property, let $\mathcal{A}_1 = (a_0, a_1, a_2, \dots)$ and $\mathcal{A}_2 = (b_0, b_1, b_2, \dots)$ be two sequences in $\Fib$. We define their sum term-wise as $\mathcal{A}_1 + \mathcal{A}_2 := (c_0, c_1, c_2, \dots)$, where $c_n = a_n + b_n$. We must verify that the sequence $(c_n)$ obeys the Fibonacci recurrence relation, namely $c_n = c_{n-1} + c_{n-2}$ for $n \ge 2$. By algebraic manipulation, we find:
\[
c_{n-1} + c_{n-2} = (a_{n-1} + b_{n-1}) + (a_{n-2} + b_{n-2}) = (a_{n-1} + a_{n-2}) + (b_{n-1} + b_{n-2})
\]
Because both $\mathcal{A}_1$ and $\mathcal{A}_2$ are Fibonacci sequences, we can substitute $a_{n-1} + a_{n-2} = a_n$ and $b_{n-1} + b_{n-2} = b_n$. Therefore,
\[
c_{n-1} + c_{n-2} = a_n + b_n = c_n
\]
This shows that $\mathcal{F}_{a_0, a_1} + \mathcal{F}_{b_0, b_1} = \mathcal{F}_{a_0 + b_0, a_1 + b_1}$, and therefore, $\Fib$ is closed under addition.

For the second property, let $\mathcal{A} = (a_0, a_1, a_2, \dots) \in \Fib$ and $\alpha \in \mathbb{R}$. We define $\alpha \mathcal{A} := (\alpha a_0, \alpha a_1, \alpha a_2, \dots)$. We need to check that $\alpha a_n = \alpha a_{n-1} + \alpha a_{n-2}$. This is obviously true because $a_n = a_{n-1} + a_{n-2}$, and we simply multiply both sides by $\alpha$. In fact, we see that $\alpha \mathcal{F}_{a_0, a_1} = \mathcal{F}_{\alpha a_0, \alpha a_1}$.
\end{proof}

\begin{corollary}
Linear combinations of Fibonacci sequences are also Fibonacci sequences. That is, for any scalars $\alpha, \beta \in \mathbb{R}$, we have:
\[
\alpha \mathcal{F}_{a_0, a_1} + \beta \mathcal{F}_{b_0, b_1} = \mathcal{F}_{\alpha a_0 + \beta b_0, \alpha a_1 + \beta b_1}
\]
\end{corollary}

\begin{remark}
As a direct consequence of this structure, in order to find a general formula for \textit{any} sequence $\mathcal{F}_{a_0, a_1}$, it is enough to find explicit formulas for the two foundational sequences $\mathcal{F}_{1,0}$ and $\mathcal{F}_{0,1}$. This is because any sequence can be decomposed as:
\[
\mathcal{F}_{a_0, a_1} = a_0 \cdot \mathcal{F}_{1,0} + a_1 \cdot \mathcal{F}_{0,1}
\]
We will see later in the course that $\Fib$ represents a $2$-dimensional vector space. Every element in this space can be expressed by a linear combination of some two base elements (but $1$ element alone is not enough). What precisely constitutes a \qt{vector space} and \qt{dimension} will be formally defined later in the course.
\end{remark}

\subsection{Deriving the Explicit Formula}

Perhaps one of the sequences in $\Fib$ matches a form we already know? 

First, we can rule out arithmetic sequences. Except for the trivial zero sequence $\mathcal{F}_{0,0}$, no Fibonacci sequence can be arithmetic. The reason is straightforward: for an arithmetic sequence, the difference between consecutive terms must be constant, meaning $a_n - a_{n-1} = d$. However, for a Fibonacci sequence, $a_n - a_{n-1} = a_{n-2}$. This difference grows and cannot remain constant. Thus, arithmetic sequences are unsuitable for this purpose.

Next, we ask: perhaps for some choice of $a_0, a_1$, the sequence $\mathcal{F}_{a_0, a_1}$ coincides with a geometric sequence? Let us try a geometric sequence of the form $\mathcal{G} = (1, q, q^2, q^3, \dots, q^n, \dots)$. 

\begin{lemma}
A geometric sequence $\mathcal{G} = (1, q, q^2, \dots)$ with $q \ne 0$ is a Fibonacci sequence if and only if $q = \frac{1 \pm \sqrt{5}}{2}$.
\end{lemma}

\begin{proof}
For $\mathcal{G}$ to be a Fibonacci sequence, it must satisfy the recurrence relation for all $n \ge 2$:
\[
q^n = q^{n-1} + q^{n-2}
\]
Since $q \ne 0$, we can divide the entire equation by $q^{n-2}$, which simplifies the condition to a quadratic equation:
\[
q^2 = q + 1 \implies q^2 - q - 1 = 0
\]
Applying the standard quadratic formula yields two possible solutions:
\[
q = \frac{1 \pm \sqrt{5}}{2}
\]
We denote these two specific roots as $\varphi := \frac{1 + \sqrt{5}}{2} \approx 1.618$ (the golden ratio) and $\psi := \frac{1 - \sqrt{5}}{2} \approx -0.618$. Consequently, we have discovered two purely geometric sequences that belong to $\Fib$:
\[
\mathcal{F}_{1, \varphi} = (1, \varphi, \varphi^2, \varphi^3, \dots, \varphi^n, \dots)
\]
\[
\mathcal{F}_{1, \psi} = (1, \psi, \psi^2, \psi^3, \dots, \psi^n, \dots)
\]
\end{proof}

\begin{theorem}[Binet's Formula]
\label{thm:binet_formula}
The explicit formula for the $n$-th term of the classical Fibonacci sequence $\mathcal{F}_{0,1}$ is given by \qt{Binet's Formula}:
\[
a_n = \frac{1}{\sqrt{5}} \left( \left(\frac{1+\sqrt{5}}{2}\right)^n - \left(\frac{1-\sqrt{5}}{2}\right)^n \right)
\]
\end{theorem}

\begin{proof}
So far, we have found formulas for two specific sequences, $\mathcal{F}_{1, \varphi}$ and $\mathcal{F}_{1, \psi}$. Our original goal is to find a formula for the classical sequence $\mathcal{F}_{0,1}$. 

We use the linearity principle established in \cref{prop:fib_closure}. We seek scalars $\alpha, \beta \in \mathbb{R}$ such that:
\[
\alpha \cdot \mathcal{F}_{1, \varphi} + \beta \cdot \mathcal{F}_{1, \psi} = \mathcal{F}_{0,1}
\]
Comparing the first two terms ($n=0$ and $n=1$) of these sequences yields a system of linear equations:
\begin{equation}
\begin{cases}
\alpha \cdot 1 + \beta \cdot 1 = 0 \\
\alpha \cdot \varphi + \beta \cdot \psi = 1
\end{cases}
\end{equation}
From the first equation, we immediately deduce that $\beta = -\alpha$. Substituting this into the second equation gives:
\[
\alpha \cdot \varphi - \alpha \cdot \psi = 1 \implies \alpha (\varphi - \psi) = 1
\]
Notice that the difference $\varphi - \psi = \sqrt{5}$. Therefore, $\alpha = 1/\sqrt{5}$ and $\beta = -1/\sqrt{5}$. This leads directly to:
\[
a_n = \frac{1}{\sqrt{5}} (\varphi^n - \psi^n) = \frac{1}{\sqrt{5}} \left( \left(\frac{1+\sqrt{5}}{2}\right)^n - \left(\frac{1-\sqrt{5}}{2}\right)^n \right)
\]
\end{proof}

\subsection{The Upshot: Asymptotics and Matrix Form}

The derivation above yields two profound observations regarding the growth and representation of the Fibonacci sequence.

\begin{importantremark}
Clearly, as $n$ grows, the term involving $\psi$ becomes negligible. Since $|\psi| = |\frac{1-\sqrt{5}}{2}| \approx 0.618 < 1$, it follows that $\psi^n \to 0$ as $n \to \infty$. Therefore, for large $n$, we have the approximation $a_n \approx \frac{1}{\sqrt{5}} \varphi^n$. In fact, because $\left| \frac{\psi^n}{\sqrt{5}} \right| < \frac{1}{2}$ for all $n \ge 0$, we arrive at the following precise conclusion:
\[
\implies a_n \text{ is the \textit{closest integer} to } \frac{1}{\sqrt{5}} \cdot \left(\frac{1+\sqrt{5}}{2}\right)^n
\]
Furthermore, this implies that the ratio of consecutive terms satisfies:
\begin{equation}
\lim_{n \to \infty} \frac{a_{n+1}}{a_n} = \varphi
\end{equation}
\end{importantremark}

Finally, we can represent the recurrence relation using the language of matrices. Defining a vector $v_n = \binom{a_{n+1}}{a_n}$, the relation $a_{n+1} = a_n + a_{n-1}$ can be written as:
\begin{equation}
\begin{pmatrix} a_{n+1} \\ a_n \end{pmatrix} = \begin{pmatrix} 1 & 1 \\ 1 & 0 \end{pmatrix} \begin{pmatrix} a_n \\ a_{n-1} \end{pmatrix} \implies v_n = \begin{pmatrix} 1 & 1 \\ 1 & 0 \end{pmatrix}^n \begin{pmatrix} a_1 \\ a_0 \end{pmatrix}
\end{equation}
This formulation shows that the behavior of the sequence is entirely governed by the powers of the matrix $M = \left(\begin{smallmatrix} 1 & 1 \\ 1 & 0 \end{smallmatrix}\right)$.

\begin{exercise}[Interesting Exercises]
\begin{enumerate}
    \item \textbf{General Formula:} Find an explicit formula for the general Fibonacci sequence $\mathcal{F}_{a_0, a_1}$ that depends only on the initial values $a_0, a_1$ and the index $n$.
    \item \textbf{The Shift Operator:} Let $S: \Fib \to \Fib$ be the shift map defined by $S(a_0, a_1, \dots) = (a_1, a_2, \dots)$. Formally verify that $S$ is a \textit{linear map}.
    \item \textbf{Golden Ratio Continued Fraction:} Consider $\varphi = \frac{1+\sqrt{5}}{2}$.
        \[
        \left.
        \begin{aligned}
        \varphi
        = 1+\cfrac{1}{\,1+\cfrac{1}{\,1+\cfrac{1}{\,1+\cfrac{1}{\,1+\cdots}}}}
        \end{aligned}
        \qquad
        \right\}
        \qquad
        \left.
        \begin{aligned}
        &\text{what is actually}\\
        &\text{the meaning of}\\
        &\text{this?}
        \end{aligned}
        \right.
        \]
    \item \textbf{Non-existence of Arithmetic Fibonacci Sequences:} Complete the details to show that no Fibonacci sequence $\mathcal{F}_{a_0, a_1}$ (except for the trivial zero sequence) can be an arithmetic sequence.
    \item \textbf{Convergence:} For the classical Fibonacci sequence $a_n = \mathcal{F}_{0,1}$, calculate the limit \[\lim_{n \to \infty} \frac{a_n}{a_{n-1}}\].
    \item \textbf{Cassini's Identity:} Show that $a_{n+1} a_{n-1} - a_n^2 = (-1)^n$ for all $n \ge 1$. Hint: Use the determinant of the matrix $M^n$.
\end{enumerate}
\end{exercise}